\section{X. 국제법 (INTERNATIONAL
LAW)}\label{x.-uxad6duxc81cuxbc95-international-law}

\subsection{1. 의문의 원천들 (SOURCES OF
DOUBT)}\label{uxc758uxbb38uxc758-uxc6d0uxcc9cuxb4e4-sources-of-doubt}

본서에서 그토록 중요한 자리를 부여받은 일차 규칙들과 이차 규칙들의
결합이라는 관념(idea)은, 법리학적 극단들 사이의 중간으로 간주될 수 있다.
법이론은 때로는 위협으로 뒷받침된 명령이라는 단순한 관념에서, 때로는
도덕(morality)이라는 복잡한 관념에서 법 이해의 열쇠를 찾으려 해왔다.
법은 이 둘 모두와 분명 많은 친연성과 연결을 가진다. 그러나 우리가
보았듯이, 이것들을 과장하고 법을 다른 사회통제 수단들과 구별하는 특수한
특징들을 흐릴 상존하는 위험이 있다. 우리가 중심적인 것으로 취한 관념의
미덕은, 법, 강제(coercion), 도덕 사이의 복수적 관계들을 그것들이 실제로
그러한 대로 보게 해 주고, 이것들이 있다면 어떤 의미에서 필연적인지를
새롭게 고려하게 해 준다는 데 있다.

일차 규칙들과 이차 규칙들의 결합이라는 관념이 이러한 미덕들을 가지고
있고, 또 이러한 특징적인 규칙 결합의 존재를 ‘법체계(legal system)’라는
표현 적용의 충분조건으로 다루는 것이 용례(usage)에 부합할 것임에도,
우리는 ‘법(law)’이라는 단어가 반드시 그 결합의 관점에서 정의되어야
한다고 주장하지 않았다. 우리가 ‘법’이나 ‘법적’ 같은 단어들의 사용을 이런
방식으로 식별하거나 규율해야 한다는 주장을 하지 않기 때문에, 이 책은
자연스럽게 이 표현들의 사용을 위한 하나의 규칙 또는 여러 규칙들을 제공할
것으로 기대될 수 있는 ‘법’의 정의가 아니라, 법 \emphksans{개념(concept)}의
해명으로 제시된다. 이 목적에 부합하게, 우리는 지난 장에서 독일
사례들에서 제기된 주장(claim), 곧 어떤 규칙들이 현존하는 일차 및 이차
규칙 체계에 속하더라도 그 도덕적 사악함(moral iniquity) 때문에 유효한
법이라는 명칭(title)을 그 규칙들로부터 보류해야 한다는 주장을
검토하였다. 결국 우리는 이 주장을 거부하였다. 그러나 그렇게 한 것은
그러한 체계에 속하는 규칙들은 반드시 ‘법’이라고 불려야 한다는 견해와 그
주장이 충돌하기 때문도 아니고, 용례의 무게와 충돌하기 때문도 아니었다.
\textbf{(p.~214)} 오히려 우리는 도덕적으로 사악한 것을 밀어냄으로써
유효한 법들의 부류를 좁히려는 시도를 비판하였다. 그 근거는 그렇게 하는
것이 이론적 탐구나 도덕적 숙고 어느 쪽도 진전시키거나 명료하게 하지
못한다는 데 있었다. 이러한 목적들을 위해서는, 많은 용례와 부합하고
규칙들이 아무리 도덕적으로 사악하더라도 그것들을 법으로 간주할 수 있게
하는 더 넓은 개념이, 검토 결과 적절한 것으로 드러났다.

국제법(international law)은 우리에게 그 반대의 사례를 제시한다. 지난
150년의 용례와 부합하게 여기서 ‘법’이라는 표현을 사용하더라도, 국제
입법부, 강제관할권을 가진 법원들, 중앙적으로 조직된 제재들의 부재는
적어도 법이론가들의 마음속에 의구심들을 불러일으켜 왔다. 이러한 제도들의
부재는 국가들에 대한 규칙들이, 오직 의무의 일차 규칙들만으로 이루어진
단순한 사회구조의 형식과 닮아 있음을 뜻한다. 개인들의 사회에서 이런
형식을 발견할 때 우리는 그것을 발전된 법체계와 대조하는 데 익숙하다.
우리가 보이겠지만, 국제법은 입법부와 법원들을 마련하는 변경과 재판의
이차 규칙들을 결여할 뿐 아니라, 법의 ‘원천들(sources)’을 특정하고 그
규칙들의 식별을 위한 일반 기준들을 제공하는 통합적 승인 규칙도
결여한다고 기실 논변할 수 있다. 이러한 차이들은 기실 두드러지며,
‘국제법은 정말로 법인가?’라는 질문은 거의 제쳐 놓을 수 없다. 그러나 이
경우에도 우리는 많은 사람들이 느끼는 의문들을 기존 용례를 단순히
환기시키는 것으로 물리치지 않을 것이며, 일차 및 이차 규칙들의 결합이
‘법체계’라는 표현을 올바르게 사용하기 위한 충분조건일 뿐 아니라
필요조건이기도 하다는 근거 위에서 그 의문들을 단순히 확증하지도 않을
것이다. 그 대신 우리는 느껴져 온 의문들의 상세한 성격을 탐구하고, 독일
사례에서처럼, ‘국제법’을 말하는 공통의 더 넓은 용례가 어떤 실천적 또는
이론적 목적을 방해할 개연성이 있는지를 물을 것이다.

우리가 국제법의 성격에 관한 이 문제에 단 하나의 장만을 할애하겠지만,
일부 저자들은 이 질문을 훨씬 더 짧게 처리할 것을 제안해 왔다. 그들에게는
‘국제법은 정말로 법인가?’라는 질문이 오직 단어들의 의미에 관한 사소한
질문을 사물들의 본성에 관한 \textbf{(p.~215)} 중대한 질문으로 오인했기
때문에 생겨났거나 살아남은 것처럼 보였다. 곧 국제법을 국내법(municipal
law)과 구별하는 사실들은 명료하고 잘 알려져 있으므로, 확정되어야 할
유일한 질문은 우리가 기존 관행을 준수할 것인지 아니면 그것에서 벗어날
것인지이며, 이것은 각자가 스스로 확정할 문제라는 것이다. 그러나 이
질문을 이렇게 짧게 처리하는 방식은 확실히 너무 짧다. 이론가들이
‘법’이라는 단어를 국제법으로 확장하는 데 주저하게 된 이유들 가운데, 같은
단어를 많은 서로 다른 것들에 적용하는 일을 무엇이 정당화하는가에 관한
지나치게 단순하고 기실 터무니없는(absurd) 견해가 어느 정도 역할을 한
것은 사실이다. 일반적인 분류 용어들의 확장을 통상적으로 인도하는 원칙
유형들의 다양상(variety)은 법리학에서 너무 자주 무시되어 왔다. 그럼에도
국제법에 관한 의문의 원천들은 단어 사용에 관한 이러한 잘못된 견해들보다
더 깊고 더 흥미롭다. 더욱이 이 질문을 짧게 처리하는 방식이 제시하는 두
대안, 곧 ‘기존 관행을 준수할 것인가, 아니면 그것에서 벗어날 것인가?’는
망라적이지 않다. 그 둘 외에도, 실제로 기존 용례를 인도해 온 원칙들을
명시하고 검토하는 대안이 있기 때문이다.

제안된 이 짧은 처리 방식은 우리가 고유명사를 다루고 있었다면 기실
적절했을 것이다. 누군가 ‘런던’이라 불리는 장소가 \emphksans{정말로}
런던인지를 묻는다면, 우리가 할 수 있는 일은 그에게 관행(convention)을
환기시키고 그가 그것을 따르거나 취향에 맞는 다른 이름을 선택하도록 남겨
두는 것뿐이다. 그런 경우에 어떤 원칙에 따라 런던이 그렇게 불리게
되었는지, 또 이 원칙이 받아들일 만한 것인지를 묻는 것은 부조리할 것이다.
이것이 부조리한 이유는, 고유명사의 배분은 \emphksans{오직} 임시적(\emph{ad
hoc}) 관행에만 기초하는 반면, 어떤 진지한 학문분야의 일반 용어들의
확장은, 그 원칙 또는 이론적 근거(rationale)가 무엇인지 분명하지 않을
수는 있어도, 결코 원칙이나 근거 없이 이루어지지 않기 때문이다. 현재
경우처럼 사실상 ‘우리는 그것이 법이라고 불린다는 것은 알지만, 그것이
정말 법인가?’라고 말하는 이들이 그 확장에 의문을 제기할 때, 요구되는
것은, 분명 흐릿하게나마, 그 원칙을 명시하고 그 자격
근거들(credentials)을 심사하라는 것이다.

우리는 국제법의 법적 성격에 관한 두 가지 주요한 의문의 원천들과,
그것들과 함께, 이론가들이 이러한 의문들에 대응하기 위해 취해 온 조치들을
검토할 것이다. \textbf{(p.~216)} 두 형태의 의문은 모두, 법이 무엇인가에
관한 명료하고 표준적인 사례로 간주되는 국내법과 국제법을 불리하게
비교하는 데서 생겨난다. 첫째 의문은 법을 근본적으로 위협으로 뒷받침된
명령들의 문제로 보는 구상(conception)에 깊이 뿌리박고 있으며, 국제법의
\emphksans{규칙들(rules)}의 성격을 국내법의 규칙들의 성격과 대조한다. 둘째
형태의 의문은, 국가들이 근본적으로 법적 의무의 수범자들(subjects)이 될
수 없다는 모호한 믿음에서 솟아나며, 국제법의 \emphksans{수범자들(subjects)}의
성격을 국내법의 수범자들의 성격과 대조한다.

\subsection{2. 의무들과 제재들 (OBLIGATIONS AND
SANCTIONS)}\label{uxc758uxbb34uxb4e4uxacfc-uxc81cuxc7acuxb4e4-obligations-and-sanctions}

우리가 검토할 의문들은 국제법 책들의 첫 장들에서 흔히 ‘국제법은 어떻게
구속력 있을 수 있는가?’라는 질문 형식으로 표현된다. 그러나 이렇게 즐겨
쓰이는 질문 형식에는 매우 혼란스러운 점이 있다. 그것을 다루기 전에
우리는 답이 결코 분명하지 않은 선행 질문에 직면해야 한다. 이 선행 질문은
이렇다. 법체계 전체에 관하여 그것이 ‘구속력 있다(binding)’고 말할 때
무엇을 뜻하는가? 한 체계의 특정 규칙이 특정 사람에게 구속력 있다고 하는
진술은 법률가들에게 익숙하고 의미상 상당히 명료하다. 우리는 그것을,
문제의 규칙이 유효한 규칙이며 그 규칙 아래에서 문제의 사람에게 어떤
의무(obligation) 또는 책무(duty)가 있다는 단언으로 바꾸어 말할 수 있다.
이 외에도 이 형식의 더 일반적인 진술들이 이루어지는 몇몇 상황들이 있다.
우리는 어떤 상황들에서 어느 법체계가 특정 사람에게 적용되는지 의문을
가질 수 있다. 그러한 의문들은 국제사법(conflict of laws)이나
공국제법(public international law)에서 생길 수 있다. 전자의 경우 우리는
특정 거래에 관하여 프랑스법과 영국법 중 어느 것이 특정 사람에게 구속력
있는지를 물을 수 있고, 후자의 경우 우리는 예컨대 적국이 점령한 벨기에의
주민들이 망명정부가 벨기에법이라고 주장한 것에 구속되었는지, 아니면
점령국의 명령들(ordinances)에 구속되었는지를 물을 수 있다. 그러나 이 두
경우 모두에서 질문들은 어떤 법체계, 곧 국내법이든 국제법이든 어떤 체계
\emphksans{내부에서(within)} 생겨나고 그 체계의 규칙들이나 원칙들을 참조하여
확정되는 법문제들(questions of law)이다. 그것들은 규칙들의 일반적 성격을
문제 삼지 않고, 다만 주어진 상황에서 특정 사람들 또는 거래들에 대한
규칙들의 범위 또는 적용가능성(applicability)을 문제 삼는다.
\textbf{(p.~217)} 분명히 ‘국제법은 구속력 있는가?’라는 질문과 그 동류
질문들인 ‘국제법은 어떻게 구속력 있을 수 있는가?’ 또는 ‘무엇이 국제법을
구속력 있게 만드는가?’는 다른 차원의 질문들이다. 그것들은 적용가능성이
아니라 국제법의 일반적 법적 지위에 관한 의문을 표현한다. 이 의문은
‘이러한 규칙들이 언제라도 의미 있게 그리고 참되게 의무들을 발생시킨다고
말해질 수 있는가?’라는 형식으로 더 솔직하게 표현될 것이다. 책들의
논의들이 보여 주듯이, 이 점에 관한 의문의 한 원천은 단순히 그 체계에
중앙적으로 조직된 제재들이 부재한다는 것이다. 이것은 국내법과의 불리한
비교 지점 하나인데, 국내법의 규칙들은 의심할 여지 없이 ‘구속력 있는’
것으로 그리고 법적 의무의 전형들로 간주된다. 이 단계에서 더 나아가는
논변은 단순하다. 이러한 이유로 국제법의 규칙들이 ‘구속력 있는’ 것이
아니라면, 그것들을 법으로 분류하는 일을 진지하게 받아들이는 것은 확실히
방어될 수 없다는 것이다. 일상 언어의 방식들이 아무리 관대하더라도,
이것은 간과하기에는 너무 큰 차이다. 법의 본성에 관한 모든
사변(speculation)은, 법의 존재가 적어도 일정한 행위를
의무적인(obligatory) 것으로 만든다는 가정에서 출발한다.

이 논변을 검토하면서 우리는 국제체계의 사실들에 관한 모든 의문점을 그
논변에 유리하게 보아 줄 것이다. 우리는 국제연맹 규약 제16조도 유엔 헌장
제7장도 국내법의 제재들과 동일시될 수 있는 어떤 것도 국제법에 도입하지
않았다고 받아들일 것이다. 한국전쟁에도 불구하고, 또 수에즈 사건에서 어떤
교훈(moral)이 도출될 수 있더라도, 우리는 헌장의 법집행
규정들(provisions)은 그 사용이 중요할 때마다 거부권에 의해 마비될
개연성이 있으며, 단지 종이 위에만 존재한다고 말해야 한다고 상정할
것이다.

국제법은 조직화된 제재들이 결여되어 있기 때문에 구속력 있지 않다고
논변하는 것은, 법이 본질적으로 위협으로 뒷받침된 명령들의 문제라는 이론
속에 들어 있는 의무 분석을 묵시적으로 받아들이는 것이다. 이 이론은,
우리가 보았듯이, ‘의무를 가진다’ 또는 ‘구속된다’를 ‘불복종에 대해 위협된
제재나 처벌을 받을 개연성이 있다’와 동일시한다. 그러나 우리가
논변했듯이, 이러한 동일시는 모든 법적 사고와 담론에서 \textbf{(p.~218)}
의무(obligation)와 책무(duty)의 관념들(ideas)이 수행하는 역할을
왜곡한다. 실효적인 조직화된 제재들이 있는 국내법에서조차, 우리는
제III장에서 제시된 다양한 이유들 때문에, ‘나는(너는) 불복종 때문에
고통을 겪을 개연성이 있다’라는 외부 예측 진술의 의미와, ‘나는(너는)
이렇게 행위할 의무를 가진다’라는 내부 규범 진술의 의미를 구별해야 한다.
후자는 특정 사람의 상황을 행위의 지도 기준들로 수용된 규칙들의 관점에서
평가한다. 모든 규칙들이 의무들 또는 책무들을 발생시키는 것은 아니라는
점은 참이다. 또한 그렇게 하는 규칙들이 일반적으로 사적 이익의 어느 정도
희생을 요구하고, 일반적으로 부합(conformity)하라는 진지한 요구들과
일탈들에 대한 집요한 비판에 의해 뒷받침된다는 점도 참이다. 그러나 우리가
예측적 분석과, 법을 본질적으로 위협으로 뒷받침된 명령으로 보는 그 모태
구상(parent conception)으로부터 벗어나고 나면, 의무라는 규범적 관념을
조직화된 제재들에 의해 뒷받침되는 규칙들로 제한할 좋은 이유는 없어
보인다.

그러나 우리는 그 논변의 또 다른 형태를 고려해야 한다. 그것은 위협된
제재들의 개연성으로 의무를 정의하는 데 묶여 있지 않기 때문에 더
그럴듯하다. 회의론자는, 우리가 스스로 강조했듯이, 국내체계 안에는
정당하게 필수적이라고 불리는 일정한 규정들(provisions)이 있다고 지적할
수 있다. 그 가운데에는 폭력의 자유로운 사용을 금지하는 의무의 일차
규칙들과, 이 규칙들 및 다른 규칙들에 대한 제재로서 힘(force)의 공식적
사용을 마련하는 규칙들이 있다. 그러한 규칙들과 그것들을 뒷받침하는
조직화된 제재들이 이 의미에서 국내법에 필수적이라면, 그것들은 국제법에도
똑같이 필수적이지 않은가? 그렇다는 주장은 이것이 ‘구속력
있는(binding)’이나 ‘의무(obligation)’ 같은 단어들의 바로 그 의미로부터
나온다고 고집하지 않고도 유지될 수 있다.

이 형태의 논변에 대한 답은, 국내법의 지속적인 심리적·물리적 배경을
이루는 인간 존재와 그 환경에 관한 그 기초적 진리들(elementary
truths)에서 발견된다. 신체적 힘과 취약성에서 근사적으로 동등한 개인들의
사회에서 물리적 제재들은 필수적이면서도 가능하다. 그것들은 법의 제약들에
자발적으로 복종하려는 이들이, 그러한 제재들이 없다면 스스로는 법을
존중하지 않으면서 타인들의 법 존중에서 오는 이익을 거둘
악행자들(malefactors)의 단순한 희생자가 되지 않게 하기 위해 요구된다.
서로 가까운 곳에서 사는 개인들 사이에서는 \textbf{(p.~219)} 공개적
공격에 의하지 않더라도 속임수로 타인들에게 손해를 가할 기회들이 매우
크고, 빠져나갈 가망(chances)도 상당하므로, 가장 단순한 사회 형태들을
제외하면 단순한 자연적 억제요인들만으로는 법에 복종하기에는 너무
사악하거나, 너무 어리석거나, 너무 약한 이들을 제어하기에 적절할 수 없다.
그러나 같은 근사적 동등성의 사실과 제약 체계에 복종할 때의 분명한 이점들
때문에, 악행자들의 어떤 결합도 그 체계의 유지에 자발적으로 협력하려는
이들의 힘을 넘어설 개연성은 없다. 국내법의 배경을 이루는 이러한
상황들에서는, 제재들이 비교적 작은 위험으로 악행자들에게 성공적으로
사용될 수 있고, 그 제재들의 위협은 존재할 수 있는 자연적 억제요인들에
많은 것을 더할 것이다. 그러나 개인들에게 성립하는 단순한 자명한 이치들이
국가들에게는 성립하지 않고, 국제법의 사실적 배경이 국내법의 사실적
배경과 매우 다르기 바로 때문에, 제재들에 대한 유사한 필수성도
없고(국제법이 그것들에 의해 뒷받침되는 것이 바람직할 수는 있더라도),
그것들을 안전하고 실효적으로 사용할 유사한 전망도 없다.

그 까닭은 국가들 사이의 침략이 개인들 사이의 침략과 매우 다르기
때문이다. 국가들 사이의 폭력 사용은 공적일 수밖에 없고, 국제 경찰력이
없더라도 그것이, 경찰력이 없는 경우의 살인이나 절도처럼, 침략자와 피해자
사이의 문제로 남으리라는 확실성은 거의 있을 수 없다. 전쟁을 개시하는
것은 최강대국에게조차 합리적 확신을 가지고 예측하기 드문 결과를 위해
많은 것을 위험에 거는 일이다. 다른 한편, 국가들의 불평등 때문에 국제질서
편에 선 이들의 결합된 힘이 침략의 유혹을 받는 세력들(powers)보다 우세할
개연성이 있다는 상시적 보장은 있을 수 없다. 따라서 제재들의 조직과
사용은 두려운 위험들을 수반할 수 있고, 그것들의 위협은 자연적
억제요인들에 거의 보태는 것이 없을 수 있다. 이렇게 매우 다른 사실의
배경에 맞서 국제법은 국내법과 다른 형식으로 발전해 왔다. 현대 국가의
인구 안에서 범죄의 조직화된 억압과 처벌이 없다면 폭력과 절도는 매시간
예상될 것이다. 그러나 국가들의 경우에는 참혹한 전쟁들 사이에 오랜 평화의
해들이 끼어들어 왔다. 이 평화의 해들은 전쟁의 위험과 이해관계 및
국가들의 상호 필요들을 고려할 때에만 \textbf{(p.~220)} 합리적으로 기대될
수 있다. 그러나 그것들은, 다른 것들 가운데 어느 중앙 기관에 의한 집행도
마련하지 않는다는 점에서 국내법의 규칙들과 다른 규칙들에 의해 규율할
가치가 있다. 그럼에도 이 규칙들이 요구하는 것은 의무적인(obligatory)
것으로 생각되고 말해진다. 규칙들에 부합(conformity)하라는 일반적 압력이
있고, 청구들(claims)과 시인들(admissions)은 그것들에 근거하며, 그것들의
위반은 보상에 대한 집요한 요구들뿐 아니라 보복조치들과 대응조치들을
정당화하는 것으로 여겨진다. 규칙들이 무시될 때, 그것은 규칙들이 구속력
있지 않다는 근거 위에서 이루어지는 것이 아니다. 오히려 사실들을 감추려는
노력들이 이루어진다. 물론 그러한 규칙들은 국가들이 싸우려 하지 않는
쟁점들에 관한 한에서만 실효적이라고 말할 수 있다. 그럴 수 있으며, 그것은
그 체계의 중요성과 인류에 대한 가치에 불리하게 반영될 수 있다. 그러나
그럼에도 이만큼이라도 확보될 수 있다는 사실은, 조직화된 제재들이
국내법에 필수적이라는 점, 곧 국내법의 물리적·심리적 사실 배경에서의 그
필수성으로부터, 매우 다른 배경의 국제법은 그것들 없이는 어떤 의무도
부과하지 않고 ‘구속력 있지’ 않으며 따라서 ‘법’이라는 명칭을 받을 가치가
없다는 결론으로 가는 단순한 추론은 할 수 없음을 보여 준다.

\subsection{3. 의무와 국가들의 주권 (OBLIGATION AND THE SOVEREIGNTY OF
STATES)}\label{uxc758uxbb34uxc640-uxad6duxac00uxb4e4uxc758-uxc8fcuxad8c-obligation-and-the-sovereignty-of-states}

영국, 벨기에, 그리스, 소비에트 러시아는 국제법 아래에서 권리들과
의무들을 가지며, 따라서 그 수범자들(subjects)에 속한다. 그것들은
일반인은 독립적이라고 생각하고 법률가는 ‘주권적(sovereign)’이라고 인정할
국가들의 임의적 사례들이다. 국제법의 의무적 성격에 관한 가장 지속적인
당혹의 원천들 가운데 하나는, 주권적인 국가가 국제법에 의해
‘구속될(bound)’ 수도 있고 국제법 아래에서 의무를 가질 수도 있다는 사실을
받아들이거나 설명하는 데서 느껴진 어려움이었다. 이러한 형태의 회의론은
어떤 의미에서는 국제법이 제재들을 결여하기 때문에 구속력 있지 않다는
반론보다 더 극단적이다. 왜냐하면 그 반론은 언젠가 국제법이 제재 체계에
의해 강화된다면 대응될 수 있지만, 현재의 반론은 동시에 주권적이고 법의
적용을 받는(subject to law) 국가라는 구상 안에 존재한다고 말해지거나
느껴지는 근본적 비일관성에 기초하기 때문이다.

이 반론의 검토는, 주권이라는 관념(notion)을 \textbf{(p.~221)} 국가
\emphksans{안의(within)} 입법부나 다른 어떤 요소 또는 인격에 적용하는 것이
아니라 국가 자체에 적용하는 경우를 면밀히 살피는 일을 포함한다.
법리학에서 ‘주권적(sovereign)’이라는 단어가 나타날 때마다, 그것을
자기보다 낮은 자들 또는 수범자들(subjects)에 대하여 자신의 말이 법인 법
위의 인격이라는 관념(idea)과 결부시키는 경향이 있다. 우리는 이 책의 앞
장들에서 이 매혹적인 관념(notion)이 국내법체계의 구조에 얼마나 나쁜
안내자인지를 보았다. 그러나 그것은 국제법 이론에서는 훨씬 더 강력한
혼란의 원천이었다. 물론 국가를 이러한 방식으로, 마치 일종의 초인, 곧
본래 법에 매이지 않지만 그 수범자들에게는 법의 원천인 존재인 것처럼
생각하는 것은 \emphksans{가능하다}. 16세기 이래 국가와 군주의 상징적
동일시(`L'état c'est moi')는 이러한 관념을 부추겼을 수 있고, 이 관념은
정치이론뿐 아니라 법이론의 많은 부분에 의심스러운 영감을 제공해 왔다.
그러나 국제법의 이해를 위해서는 이러한 연상들을 떨쳐 내는 것이 중요하다.
‘국가’라는 표현은 본래 또는 ‘본성상’ 법 바깥에 있는 어떤 인격이나 사물의
이름이 아니다. 그것은 두 사실들을 지칭하는 방식이다. 첫째, 한 영토에
거주하는 인구가 입법부, 법원들, 일차 규칙들이라는 특징적 구조를 가진
법체계가 마련한 질서 있는 정부 형식 아래에서 산다는 사실, 둘째, 그
정부가 모호하게 규정된 정도의 독립성을 누린다는 사실이다.

‘국가(state)’라는 단어가 분명 그 자체의 넓은 모호성 영역을 가지지만,
지금까지 말한 것으로 그 중심 의미를 드러내기에는 충분할 것이다. 다시
임의적 사례들을 들면, 영국이나 브라질, 미국이나 이탈리아 같은 국가들은
자기 국경 밖의 어떤 권위들이나 인격들에 의한 법적 통제와 사실적 통제
양자 모두로부터 매우 큰 정도의 독립성을 가지며, 국제법에서
‘주권국가들’로 자리매겨질 것이다. 다른 한편, 미국의 주들처럼 연방연합의
구성원인 개별 주들은 여러 다른 방식으로 연방정부와 연방헌법의 권위와
통제에 종속된다. 그러나 이 연방 주들조차 보유하는 독립성은, 이를테면
‘국가’라는 단어가 전혀 사용되지 않을 잉글랜드 카운티의 지위와 비교하면
크다. 카운티는 자기 구역을 위해 입법부의 기능들 일부를 수행하는
지방의회를 가질 수 있지만, 그 빈약한 권한들은 의회의 권한들에
\textbf{(p.~222)} 종속되며, 어떤 사소한 측면들을 제외하면 그 카운티의
구역은 나라의 나머지 부분과 같은 법들 및 정부의 적용을 받는다.

이 극단들 사이에는 질서 있는 정부를 가진 영토 단위들 사이의 종속성,
따라서 독립성의 많은 서로 다른 유형들과 정도들이 있다. 식민지, 보호령,
종주권관계(suzerainties), 신탁통치지역, 국가연합들은 이 관점에서
흥미로운 분류 문제들을 제시한다. 대부분의 경우 한 단위가 다른 단위에
종속되어 있다는 점은 법적 형식들로 표현되므로, 종속 단위의 영토 안에서
법인 것은 적어도 일정한 쟁점들에 관해서는 궁극적으로 다른 단위의 법제정
작용들(law-making operations)에 의존할 것이다.

그러나 어떤 경우에는 종속된 영토의 법체계가 그 종속성을 반영하지 않을 수
있다. 이는 그 영토가 단지 형식적으로 독립되어 있고 실제로는 꼭두각시들을
통해 외부에서 통치되기 때문일 수도 있고, 또는 종속된 영토가 내부
사무들에 대해서는 실질적 자율성을 가지지만 외부 사무들에 대해서는 그렇지
않으며, 외부 사무들에서 다른 나라에 대한 그 종속성이 그 국내법의 일부로
표현될 필요가 없기 때문일 수도 있다. 그러나 이러한 여러 방식으로 한 영토
단위가 다른 영토 단위에 종속되는 것이 그 독립성이 제한될 수 있는 유일한
형식은 아니다. 제한 요인은 그러한 다른 단위의 권력(power)이나
권위(authority)가 아니라, 서로 똑같이 독립적인 단위들에 영향을 미치는
국제적 권위일 수 있다. 국제적 권위의 많은 서로 다른 형식들과 그에
상응하는 국가 독립성의 많은 서로 다른 제한들을 상상할 수 있다. 그
가능성들에는, 많은 다른 것들 가운데, 모든 단위의 내부/외부 사무들을
규율할 법적으로 무제한적인 권한들(powers)을 가진, 영국 의회를 모델로 한
세계 입법부, 특정 사안들에 대해서만 법적 권능(legal competence)을
가지거나 구성 단위들의 특정 권리 보장에 의해 제한되는, 미국 의회를
모델로 한 연방 입법부, 법적 통제의 유일한 형식이 모두에게 적용가능한
것으로 일반적으로 수용된 규칙들로 이루어지는 체제, 마지막으로 인정되는
의무의 유일한 형식이 계약적 또는 자기부과적인 것이어서 국가의 독립성이
오직 자신의 행위에 의해서만 법적으로 제한되는 체제가 포함된다.

이러한 가능성들의 범위를 고려하는 것은 유익하다. 종속성과 독립성의
가능한 형식들과 \textbf{(p.~223)} 정도들이 많다는 점을 깨닫기만 해도,
국가들이 주권적이기 때문에 국제법의 적용을 받거나 국제법에 의해 구속될
수 ‘\emphksans{없다}’거나 어떤 특정 형식의 국제법에 의해서만 구속될 수
‘\emphksans{있다}’는 주장(claim)에 답하는 한 걸음이 되기 때문이다. 여기서
‘주권적’이라는 단어는 ‘독립적’이라는 말 이상을 뜻하지 않으며, 후자처럼
그 의미상 힘(force)이 부정적이다. 주권국가는 일정 유형들의 통제에
종속되지 \emphksans{않는} 국가이고, 그 주권은 그 국가가 자율적인 행위
영역이다. 우리가 보았듯이, 국가라는 단어의 바로 그 의미가 어느 정도의
자율성을 불러들이지만, 이것이 무제한적이어야 ‘\emphksans{한다}’거나 일정
유형들의 의무에 의해서만 제한될 수 ‘\emphksans{있다}’는 논지(contention)는
기껏해야 국가들이 여타 모든 제약들로부터 자유로워야 한다는 주장(claim)의
단언이고, 최악의 경우에는 이유 없는 독단이다. 실제로 국가들 사이에 어떤
주어진 형식의 국제적 권위가 존재한다는 것을 우리가 발견한다면, 국가들의
주권은 그만큼 제한되며, 그것은 규칙들이 허용하는 바로 그 정도를 가진다.
따라서 어떤 국가들이 주권적인지, 그리고 그 주권의 범위가 무엇인지는
규칙들이 무엇인지 알 때에만 알 수 있다. 이는 영국인이나 미국인이
자유로운지 그리고 그의 자유의 범위가 무엇인지가 영국법이나 미국법이
무엇인지 알 때에만 알 수 있는 것과 같다. 국제법의 규칙들은 많은 점에서
기실 모호하고 충돌하므로, 국가들에게 남겨진 독립성의 영역에 관한 의문은
국내법 아래 시민의 자유 범위에 관한 의문보다 훨씬 더 크다. 그럼에도
이러한 어려움들은, 국제법을 참조하지 않고 국가들에게 속한다고 가정된
절대적 주권으로부터 국제법의 일반적 성격을 연역하려는 \emphksans{선험적(a
priori)} 논변을 타당화하지 않는다.

주권 관념의 무비판적 사용이 국내법 이론과 국제법 이론 양쪽에 유사한
혼란을 퍼뜨려 왔고, 양쪽에서 유사한 교정을 요구한다는 점은 관찰할 가치가
있다. 그 영향 아래에서 우리는 모든 국내법체계에는 법적 제한들에 종속되지
않는 주권 입법자가 반드시 있어야 \emphksans{한다고} 믿게 된다. 마찬가지로
우리는 국가들이 주권적이고 스스로에 의한 제한을 제외하고는 법적 제한을
받을 수 없기 때문에 국제법은 반드시 어떤 성격을 가져야 \emphksans{한다고}
믿게 된다. 두 경우 모두, 법적으로 무제한적인 주권자가 필연적으로
존재한다는 믿음은 실제 \textbf{(p.~224)} 규칙들을 검토할 때에만 답할 수
있는 질문을 선취 판단한다. 국내법에 관한 질문은 이것이다. 이 체계에서
인정되는 최고 입법 권위(supreme legislative authority)의 범위는
무엇인가? 국제법에 관한 질문은 이것이다. 규칙들이 국가들에게 허용하는
최대 자율성 영역은 무엇인가?

따라서 현재의 반론에 대한 가장 단순한 답은, 그것이 질문들이 고려되어야
할 순서를 뒤집는다는 것이다. 국제법의 형식들이 무엇인지, 그리고 그것들이
단지 공허한 형식들인지 아닌지를 알기 전까지는 국가들이 어떤 주권을
가지는지 알 길이 없다. 많은 법리학적 논쟁은 이 원칙이 무시되었기 때문에
혼란스러워졌고, 이 원칙에 비추어 ‘의사주의(voluntarist)’ 이론들 또는
‘자기제한(auto-limitation)’ 이론들로 알려진 국제법 이론들을 고려하는
것은 유익하다. 이러한 이론들은 국가들의 (절대적) 주권을 국제법의 구속력
있는 규칙들의 존재와 조화시키려 하면서, 모든 국제적 의무들을 약속에서
생겨나는 의무처럼 자기부과적인 것으로 다루었다. 그러한 이론들은 사실
국제법에서 정치학의 사회계약론들에 대응하는 것이다. 후자는, ‘자연적으로’
자유롭고 독립적인 개인들이 그럼에도 국내법에 의해 구속된다는 사실들을,
법에 복종할 의무를 구속되는 이들이 서로, 그리고 몇몇 경우에는 자기
통치자들과 맺은 계약에서 생겨나는 의무로 다루어 설명하고자 했다. 우리는
여기서 이 이론을 문자 그대로 취할 때의 잘 알려진 반론들이나, 단지 해명적
유비로 취할 때의 그 가치에 대해서는 고려하지 않을 것이다. 그 대신 우리는
그 역사에서 의사주의 국제법 이론들에 반대하는 세 갈래 논변을 끌어낼
것이다.

첫째, 이러한 이론들은 국가들이 자기부과적 의무들에 의해서만 구속될 수
‘\emphksans{있다}’는 점이 어떻게 알려지는지, 또는 국제법의 실제 성격에 대한
어떤 검토보다 앞서 왜 그들의 주권에 관한 이러한 견해가 수용되어야
하는지를 전혀 설명하지 못한다. 그것을 뒷받침할 것이, 그것이 자주
반복되었다는 사실 외에 더 있는가? 둘째, 국가들이 그 주권 때문에
스스로에게 부과한 규칙들의 적용을 받거나 그 규칙들에 의해 구속될 수 있을
\emphksans{뿐}임을 보이도록 설계된 논변에는 비일관적인 점이 있다. ‘자기제한’
이론의 몇몇 매우 극단적인 형식들에서는 한 국가의 합의나 조약상 약정들이
그 국가가 제안하는 장래 행위에 관한 단순한 선언들로 취급되고, 이행하지
않음은 어떤 의무의 위반으로도 간주되지 않는다. \textbf{(p.~225)} 이것은
사실들과 매우 어긋나지만, 적어도 일관성이라는 이점(merit)을 가진다.
그것은 국가들의 절대적 주권이 어떤 종류의 의무와도 비일관적이어서,
의회처럼 국가도 자신을 구속할 수 없다는 단순한 이론이다. 그러나 국가가
약속, 합의 또는 조약에 의해 스스로에게 의무들을 부과할 수 있다는 덜
극단적인 견해는, 국가들이 이렇게 스스로에게 부과한 규칙들의 적용을
받거나 그 규칙들에 의해 구속될 수 있을 뿐이라는 이론과 비일관적이다.
말해진 것이든 쓰인 것이든 말들이 일정한 상황에서 약속, 합의 또는
조약으로 기능하여 의무들을 발생시키고 타인들이 청구할 수 있는 권리들을
부여하려면, 한 국가가 적절한 말들로 하겠다고 맡은 것은 무엇이든 하도록
구속된다고 규정하는 \emphksans{규칙들}이 이미 존재해야 하기 때문이다.
자기부과적 의무라는 바로 그 관념(notion) 안에 전제되어 있는 그러한
규칙들은, 그것들에 복종해야 한다는 자기부과적 의무에서 \emphksans{그것들의}
의무적 지위(obligatory status)를 도출할 수 없음이 분명하다.

주어진 국가가 해야 하도록 구속된 모든 특정 \emphksans{행위(action)}가 이론상
약속에서 그 의무적 성격을 도출할 수 있다는 것은 참이다. 그럼에도 이는
약속 등이 의무들을 창설한다는 \emphksans{규칙}이 어떤 약속과도 독립적으로 그
국가에 적용되는 경우에만 그럴 수 있다. 개인들로 이루어졌든 국가들로
이루어졌든 어떤 사회에서든, 약속, 합의 또는 조약의 말들이 의무들을
발생시키기 위해 필요하고 충분한 것은, 이를 마련하고 이러한 자기구속
작용들(self-binding operations)을 위한 절차를 특정하는 규칙들이,
보편적으로 그럴 필요는 없지만 일반적으로 인정되는 것이다. 그것들이
인정되는 곳에서는 이러한 절차들을 알고서 사용하는 개인이나 국가는,
구속되기를 선택하든 선택하지 않든 그 절차들에 의해 구속된다. 그러므로
사회적 의무의 이 가장 자발적인 형식조차, 그 규칙들에 의해 구속되는
당사자의 선택과 독립적으로 구속력 있는 몇몇 규칙들을 포함하며, 이는
국가들의 경우 그들의 주권이 그러한 모든 규칙들로부터의 자유를 요구한다는
상정과 비일관적이다.

셋째, 사실들이 있다. 우리는 방금 비판한 \emphksans{선험적} 주장, 곧 국가들은
자기부과적 의무들에 의해서만 구속될 수 ‘\emphksans{있다}’는 주장과, 국가들이
다른 체계 아래에서는 다른 방식들로 구속될 수 있었더라도 사실상 현행
국제법 규칙들 아래에서는 국가들을 위한 다른 어떤 의무 형식도 존재하지
않는다는 주장을 구별해야 한다. 물론 그 \textbf{(p.~226)} 체계가 이
완전히 합의적인(consensual) 형식의 것일 수도 있으며, 그 성격에 관한 이
견해의 단언들과 부인들은 법리학자들의 저술, 판사들의 의견, 심지어
국제법원들의 의견, 그리고 국가들의 선언들에서 발견된다. 오직 국가들의
실제 실천에 대한 냉정한 조사만이 이 견해가 옳은지 아닌지를 보여 줄 수
있다. 현대 국제법이 매우 큰 부분에서 조약법이라는 점은 참이며, 국가들의
사전 동의 없이도 국가들에게 구속력 있는 것으로 보이는 규칙들이 실제로는
동의에 기초해 있다는 것을 보이려는 정교한 시도들이 이루어져 왔다. 다만
이 동의는 오직 ‘묵시적으로(tacitly)’ 주어진 것이거나 ‘추론되어야’ 하는
것일 수 있다. 모두가 허구인 것은 아니지만, 국제적 의무의 형식들을 하나로
환원하려는 이 시도들 중 적어도 일부는, 우리가 보았듯이 국내법의
유사하지만 더 명백히 허위적인 단순화를 수행하도록 설계된 ‘묵시적
명령’이라는 관념(notion)과 같은 의심을 불러일으킨다.

모든 국제적 의무가 구속되는 당사자의 동의에서 생겨난다는 주장에 대한
상세한 검토는 여기서 수행할 수 없지만, 이 교설에 대한 두 가지 명료하고
중요한 예외들은 주목되어야 한다. 첫째는 신생국가의 경우이다. 1932년의
이라크와 1948년의 이스라엘처럼 새롭고 독립적인 국가가 존재로 출현할 때,
그 국가는 그중에서도 조약들에 구속력을 부여하는 규칙들을 포함한 국제법의
일반 의무들에 구속된다는 점은 결코 의문시된 적이 없다. 여기서 신생국가의
국제적 의무들을 ‘묵시적’ 또는 ‘추론된’ 동의에 기초시키려는 시도는 완전히
빈약해 보인다. 둘째는 국가가 영토를 획득하거나 어떤 다른 변화를 겪고, 그
변화가 이전에는 그 국가가 준수하거나 위반할 기회도 없었고 동의를 주거나
보류할 계기도 없었던 규칙들 아래에서 처음으로 의무들의
귀속양상(incidence)을 가져오는 경우이다. 이전에 바다에 접근할 수 없던
국가가 해양 영토를 획득한다면, 이것이 그 국가를 영해와 공해에 관한
국제법의 모든 규칙들의 수범자로 만들기에 충분하다는 점은 분명하다. 이
밖에도 주로 일반 조약이나 다자 조약들이 비당사자들에게 미치는 효과와
관련된 더 논쟁적인 사례들이 있다. 그러나 이 두 중요한 예외들만으로도,
모든 국제적 의무가 자기부과적이라는 일반 이론이 사실들에 대한 존중은
너무 적고 추상적 독단에 지나치게 이끌려 왔다는 의심은 정당화하기에
충분하다.

\subsection{4. 국제법과 도덕 (INTERNATIONAL LAW AND
MORALITY)}\label{uxad6duxc81cuxbc95uxacfc-uxb3c4uxb355-international-law-and-morality}

\textbf{(p.~227)} 제V장에서 우리는 의무의 일차 규칙들만으로 이루어진
단순한 사회구조 형식을 검토하였고, 가장 작고 가장 긴밀하게 결속되어
있으며 고립된 사회들을 제외하면 그것이 중대한 결함들을 겪는다는 점을
보았다. 그러한 체제는 정태적이어야 하고, 그 규칙들은 성장과 쇠퇴의 느린
과정들에 의해서만 변경된다. 규칙들의 식별은 불확실해야 하고, 특정
사례들에서 그 규칙들의 위반 사실을 확인하는 일과 위반자들에게 사회적
압력을 적용하는 일은 우발적이고, 시간을 낭비하며, 약할 수밖에 없다.
우리는 국내법에 특징적인 승인, 변경, 재판의 이차 규칙들을 이러한 서로
다른 결함들에 대한 서로 다르지만 관련된 치유책들(remedies)로 파악하는
것이 해명적임을 발견하였다. 국제법은 형식상 그러한 일차 규칙들의 체제를
닮아 있다. 그럼에도 그 흔히 정교한 규칙들의 내용은 원시 사회의 규칙들과
매우 다르고, 그 개념들, 방법들, 기법들 가운데 많은 것은 현대 국내법의
그것들과 같다. 법리학자들은 매우 자주 국제법과 국내법 사이의 이러한
형식적 차이들이 전자를 ‘도덕(morality)’으로 분류함으로써 가장 잘 표현될
수 있다고 생각해 왔다. 그러나 이런 방식으로 차이를 표시하는 것은 혼란을
불러들이는 것으로 보임이 분명하다.

때로는 국가들 사이의 관계들을 규율하는 규칙들이 오직 도덕 규칙들일
뿐이라는 고집이, 위협으로 뒷받침된 명령들로 환원될 수 없는 모든 사회구조
형식은 오직 ‘도덕’의 한 형식일 수밖에 없다는 오래된 독단주의에서 영감을
받는다. 물론 ‘도덕’이라는 단어를 이처럼 매우 포괄적인 방식으로 사용하는
것은 가능하다. 그렇게 사용되면 그것은 게임, 클럽, 예절의 규칙들, 헌법과
국제법의 기본 규정들(provisions), 그리고 잔혹함, 부정직, 거짓말의 공통된
금지처럼 우리가 통상 도덕적인 것들로 생각하는 규칙들과 원칙들이 함께
들어갈 개념적 폐지통을 제공한다. 이 절차에 대한 반론은, 이렇게
‘도덕’으로 함께 분류된 것들 사이에는 형식과 사회적 기능 양쪽에서 너무
중요한 차이들이 있어서, 그토록 조악한 분류로는 생각할 수 있는 어떤
실천적 또는 이론적 목적도 봉사될 수 없다는 것이다. 이렇게 인위적으로
넓혀진 도덕의 범주 안에서는, 그것이 흐리는 오래된 구별들을 새로 다시
획정해야 할 것이다.

\textbf{(p.~228)} 국제법의 특수한 경우에는 그 규칙들을 ‘도덕’으로
분류하는 데 저항할 여러 서로 다른 이유들이 있다. 첫째는 국가들이 종종
부도덕한 행위에 대해 서로를 비난하거나, 국제 도덕의 기준에 부응했다는
이유로 자신들이나 타국들을 칭찬한다는 점이다. 국가들이 보일 수도 있고
보이지 못할 수도 있는 미덕들 가운데 \emphksans{하나}가 국제법을 준수하는
것임은 의심할 여지가 없지만, 그것이 그 법이 도덕이라는 뜻은 아니다.
실제로 도덕의 관점에서 국가들의 행위를 평가하는 일은, 국제법 규칙들
아래에서 청구들, 요구들, 권리들과 의무들의 인정들을 정식화하는 일과 식별
가능하게 다르다. 제IX장에서 우리는 사회적 도덕의 규정적 특징들(defining
characteristics)로 취해질 수 있는 일정한 특징들을 열거하였다. 그
가운데에는 도덕 규칙들이 일차적으로 뒷받침되는 도덕적 압력의 독특한
형식이 있었다. 이것은 두려움에 대한 호소나 보복의 위협이나 보상 요구로
이루어지지 않고, 말이 건네진 사람이 문제된 도덕 원칙을 일단 상기하게
되면 죄책감이나 수치심에 이끌려 그것을 존중하고 잘못을 바로잡게(make
amends) 될 수 있으리라는 기대 속에서 이루어지는 양심에 대한 호소로
이루어진다.

국제법 아래에서의 청구들은 이러한 용어들로 표현되지 않는다. 물론
국내법에서처럼 그것들은 도덕적 호소와 결합될 수 있다. 국가들이 국제법의
다툼 있는 사안들에 관하여 서로에게 제시하는, 흔히 기술적인 논변들에서
우세한 것은 선례들, 조약들, 법리학 저술들에 대한 참조들이다. 흔히 도덕적
옳음이나 그름, 좋음이나 나쁨에 대한 언급은 전혀 이루어지지 않는다.
따라서 베이징 정부가 포모사에서 국민당 세력들을 추방할 국제법 아래의
권리를 가지고 있다거나 가지고 있지 않다는 주장은, 이것이 공정한지,
정의로운지, 또는 도덕적으로 좋은 일인지 나쁜 일인지의 질문과 매우
다르며, 특징적으로 서로 다른 논변들에 의해 뒷받침된다. 의심할 여지 없이
국가들 사이의 관계들에는, 사적 생활에서 인정되는 예의와 공손의 기준들과
유비적인, 분명히 법인 것과 분명히 도덕인 것 사이의 중간지대들이 있다.
개인적 사용을 위한 물품을 무관세로 받을 수 있도록 외교사절에게 확대된
특권에서 예시되는 국제 예양(international comity)의 영역이 그러하다.

더 중요한 구별 근거는 다음과 같다. 국제법의 규칙들은 국내법의 규칙들처럼
\textbf{(p.~229)} 흔히 도덕적으로 상당히 무관하다. 하나의 규칙은 그것이
다루는 주제들에 관하여 어떤 명료하고 고정된 규칙을 가지는 것이
편리하거나 필수적이기 때문에 존재할 수 있지, 그 특정 규칙에 어떤 도덕적
중요성이 부착되기 때문에 존재하는 것은 아닐 수 있다. 그것은 가능했던
많은 규칙들 중 하나일 수 있고, 그중 어느 하나라도 똑같이 잘 기능했을 수
있다. 따라서 국내적·국제적 법규칙들은 흔히 많은 구체적 세부사항을 담고,
도덕 규칙들이나 원칙들 안의 요소들로서는 이해될 수 없을 자의적 구분들을
긋는다. 사회적 도덕의 가능한 내용에 관하여 우리가 독단적이어서는 안
된다는 점은 참이다. 제IX장에서 보았듯이 한 사회 집단의 도덕은 현대
지식의 빛에서 보면 터무니없거나(absurd) 미신적인 것으로 보일 수 있는
많은 명령을 포함할 수 있다. 그러므로 우리와 매우 다른 일반적 믿음을 가진
사람들이 도로의 오른쪽이 아니라 왼쪽으로 운전하는 데 \emphksans{도덕적}
중요성을 부착하게 되거나, 두 명의 증인이 지켜본 약속을 어기면 도덕적
죄책감을 느끼지만 한 명이 지켜본 경우에는 그러한 죄책감을 느끼지 않게
된다고 상상하는 것은 어렵지만 가능하다. 그러한 낯선 도덕들이
가능하더라도, 도덕은 그것을 따르는 이들이 일반적으로 대안들보다 결코 더
선호할 만하지 않고 아무 내재적 중요성도 없다고 여기는 규칙들을
논리적으로 포함할 수 없다는 점은 여전히 참이다. 그러나 법은 도덕적으로
중요한 많은 것을 포함하기도 하지만, 바로 그러한 규칙들을 포함할 수 있고
실제로 포함한다. 따라서 도덕의 일부로 이해하기는 가장 어려울 자의적
구분들, 형식절차들(formalities), 고도로 구체적인 세부사항은 법의
자연스럽고 쉽게 이해되는 특징들이다. 도덕과 달리 법의 전형적 기능 가운데
하나는 확실성과 예측가능성을 극대화하고 청구들의 증명 또는 평가를
용이하게 하기 위해 바로 이러한 요소들을 도입하는 것이기 때문이다.
형식들과 세부사항에 대한 지나친 고려는 법으로 하여금 ‘형식주의’와
‘법률주의’라는 비난을 얻게 했다. 그러나 이러한 악덕들이 법의 독특한
성질들 일부를 과장한 것임을 기억하는 것이 중요하다.

이 이유 때문에, 유효하게 작성된 유언장이 몇 명의 증인을 가져야 하는지를
도덕이 아니라 국내법체계가 말해 주기를 기대하듯이, 교전국 선박이 급유나
수리를 위해 \textbf{(p.~230)} 중립항에 머물 수 있는 날짜 수, 영해의 폭,
그것을 측정하는 데 사용될 방법들과 같은 것들은 도덕이 아니라 국제법이
말해 주기를 기대한다. 이 모든 것들은 \emphksans{법규칙들}이 정하는 것이
필수적이고 바람직한 규정 사항들(provisions)이다. 그러나 그러한 규칙들이
몇몇 형식들 중 어느 것을 취해도 똑같이 괜찮을 수 있거나, 특정 목적들에
이르는 많은 가능한 수단들 가운데 하나로서만 중요하다는 감각이 유지되는
한, 그것들은 개인적 또는 사회적 삶에서 도덕에 특징적인 지위를 가지는
규칙들과 구별되어 남아 있다. 물론 국제법의 모든 규칙들이 이런
형식적이거나 자의적이거나 도덕적으로 중립적인 종류인 것은 아니다. 요점은
오직 법규칙들은 이런 종류일 \emphksans{수} 있고 도덕 규칙들은 그럴 \emphksans{수
없다}는 것이다.

국제법과 우리가 자연스럽게 도덕으로 생각하는 것 사이의 성격 차이에는 또
다른 측면이 있다. 일정한 실천들을 요구하거나 금지하는 법의 효과가
궁극적으로 한 집단의 도덕에 변화를 일으키는 것일 수는 있지만, 도덕
규칙들을 만들거나 폐지하는 입법부라는 관념(notion)은, 제VII장에서
보았듯이, 부조리한 것이다. 입법부는 자기 \emph{fiat}에 의해 새 규칙을
도입하고 그것에 도덕 규칙의 지위를 부여할 수 없다. 같은 수단으로 한
규칙에 전통의 지위를 부여할 수 없는 것과 마찬가지이다. 다만 왜
그러한지는 두 경우에서 같지 않을 수 있다. 따라서 도덕은 단지 입법부를
결여하거나 우연히 입법부를 가지지 않는 것이 아니다. 인간의 입법적
\emph{fiat}에 의한 변경이라는 바로 그 관념이 도덕이라는 관념에 배치된다.
우리가 도덕을 인간 행위들, 입법적 행위들이든 그 밖의 행위들이든 그것들을
평가하는 궁극적 기준으로 파악하기 때문에 그러하다. 국제법과의 대조는
명료하다. 국제법의 본성이나 기능 안에는 그 규칙들이 입법적 변경의 대상이
될 수 있다는 관념과 유사하게 비일관적인 것이 없다. 입법부의 결여는 많은
이들이 언젠가 수선되어야 할 결함으로 생각하는 바로 그런 결여일 뿐이다.

끝으로 우리는 국제법 이론에서, 제IX장에서 비판된 논변, 곧 국내법의 특정
규칙들이 도덕과 충돌할 수 있더라도 체계 전체는 그 규칙들에 복종할 도덕적
의무가 있다는 일반적으로 확산된 확신에 기초해야 하며, 이것은 특별한
예외적 사례들에서만 압도될 수 있다는 논변과의 병행성을 주목해야 한다.
국제법의 ‘토대들(foundations)’에 관한 논의에서, \textbf{(p.~231)}
최종적으로 국제법의 규칙들은 그것들에 복종할 도덕적 의무가 있다는
국가들의 확신에 기초해야 \emphksans{한다}고 자주 말해져 왔다. 그러나 이것이
국가들이 인정하는 의무들이 공적으로 조직된 제재들에 의해 집행가능하지
않다는 것 이상을 뜻한다면, 그것을 받아들일 이유는 없어 보인다. 물론 한
국가가 국제법에 의해 요구되는 어떤 행위 과정을 도덕적으로 의무적인
것으로 여겼고 그 이유로 행위했다고 우리가 말하는 것을 확실히 정당화할
상황들을 생각할 수 있다. 예컨대 조약들에 대한 신뢰가 심하게 흔들리면
인류에게 명백한 해악이 따를 것이기 때문에, 또는 자신이 그 차례에
타인들에게 부담이 떨어졌을 때 과거에 이익을 얻었던 규약(code)의 성가신
부담들을 짊어지는 것이 공정할 뿐이라는 감각 때문에, 한 국가는 부담스러운
조약의 의무들을 계속 이행할 수 있다. 도덕적 확신에 관한 그러한
사안들에서 정확히 누구의 동기들, 생각들(thoughts), 감정들이 국가에
귀속되어야 하는지는 여기서 우리를 붙들어 둘 필요가 없는 질문이다.

그러나 그러한 도덕적 의무의 감각이 있을 \emphksans{수} 있다 하더라도,
국제법의 존재 조건으로 그것이 왜 또는 어떤 의미에서 반드시 존재해야
\emphksans{하는지} 보기는 어렵다. 국가들의 실천에서 일정한 규칙들이 일정한
희생을 대가로 하면서도 반복적으로 일양적인 존중을 받는다는 점, 그것들을
참조하여 청구들이 정식화된다는 점, 규칙들의 위반은 위반자를 중대한
비판에 노출시키고 보상 또는 보복에 대한 청구들을 정당화하는 것으로
여겨진다는 점은 분명하다. 이것들은 확실히 국가들 사이에 그들에게
의무들을 부과하는 규칙들이 존재한다는 진술을 뒷받침하는 데 요구되는 모든
요소들이다. 어떤 사회에서 ‘구속력 있는’ 규칙들이 존재한다는 입증은,
단순히 그것들이 그러한 것으로 생각되고, 말해지고, 기능한다는 것이다.
‘토대들’의 방식으로 무엇이 더 요구되며, 더 요구된다면 왜 그것은 도덕적
의무의 토대여야 하는가? 물론 압도적 다수가 규칙들을 수용하고 그것들의
유지에 자발적으로 협력하지 않는 한, 규칙들은 국가들 사이의 관계에서
존재하거나 기능할 수 없다는 것은 참이다. 또한 규칙들을 깨뜨리거나
깨뜨리겠다고 위협하는 이들에게 행사되는 압력은 흔히 비교적 약하고,
대체로 탈중심화되어 있거나 조직화되어 있지 않다는 점도 참이다. 그러나
훨씬 더 강하게 강제적인 국내법체계를 자발적으로 수용하는 개인들의 경우와
마찬가지로, 그러한 체계를 자발적으로 지지하는 동기들은 극도로 다양할 수
있다. 어떤 \textbf{(p.~232)} 형식의 법질서든 그것에 부합하는 것이
도덕적으로 의무적이라는 감각이 일반적으로 퍼져 있을 때 가장 건강할 수는
있다. 그럼에도 법에 대한 고수(adherence)는 그것에 의해 동기부여되지 않을
수 있고, 장기적 이익의 계산이나 전통을 계속하려는 바람, 또는 타인들에
대한 사심 없는 관심에 의해 동기부여될 수 있다. 개인들 사이에서든 국가들
사이에서든, 이들 중 어느 하나를 법의 존재에 필요한 조건으로 식별할 좋은
이유는 없어 보인다.

\subsection{5. 형식과 내용의 유비들 (ANALOGIES OF FORM AND
CONTENT)}\label{uxd615uxc2dduxacfc-uxb0b4uxc6a9uxc758-uxc720uxbe44uxb4e4-analogies-of-form-and-content}

순진한 눈에는, 입법부, 강제관할권을 가진 법원들, 공적으로 조직된
제재들을 결여한 국제법의 형식적 구조가 국내법의 그것과 매우 달라 보인다.
그것은 우리가 말했듯이 내용에서는 전혀 아니지만 형식에서는 일차법 또는
관습법의 단순한 체제를 닮아 있다. 그러나 어떤 이론가들은 국제법이
‘법’이라고 불릴 자격(title)을 회의론자에 맞서 방어하려는 초조함 속에서,
이러한 형식적 차이들을 최소화하고 국제법 안에서 입법이나 국내법의 다른
바람직한 형식적 특징들과 관련하여 발견될 수 있는 유비들을 과장하려는
유혹에 굴복하였다. 그래서 패전국이 영토를 양도하거나, 의무들을
부담하거나, 어떤 축소된 독립 형식을 받아들이는 조약으로 끝나는 전쟁은
본질적으로 입법적 행위라고 주장되어 왔다. 입법처럼 그것은 부과된 법적
변경이기 때문이다. 오늘날 이 유비에 감명받거나, 그것이 국제법도 국내법과
동등하게 ‘법’이라고 불릴 자격(title)을 가진다는 점을 보이는 데 도움이
된다고 생각할 사람은 거의 없을 것이다. 국내법과 국제법 사이의 두드러진
차이들 가운데 하나는, 전자는 보통 폭력으로 강취된 합의들의 유효성을
인정하지 않지만 후자는 인정한다는 점이기 때문이다.

‘법’이라고 불릴 자격(title)이 그것들에 달려 있다고 생각하는 이들은 여러
다른, 더 존중할 만한 유비들을 강조해 왔다. 국제사법재판소와 그 전신인
상설국제사법재판소의 판결들이 거의 모든 사례에서 당사자들에 의해 제대로
이행되어 왔다는 사실은, 국내법 법원들과 대조적으로 어떤 국가도 그 사전
동의 없이는 이러한 국제재판소들 앞에 세워질 수 없다는 사실을 그것이
어떻게든 상쇄하는 것처럼 자주 강조되어 왔다. 또한 국내법에서 제재로서
법적으로 규율되고 공적으로 집행되는 \textbf{(p.~233)} 힘의 사용과
‘탈중심화된 제재들’, 곧 국제법 아래의 자기 권리들이 타국에 의해
침해되었다고 주장하는 국가가 전쟁이나 무력 보복에 의존하는 것 사이에서도
유비들이 발견되어 왔다. 어느 정도 유비가 있다는 것은 명백하다. 그러나 그
의의는, 국내법 법원은 ‘자력구제’의 옳음과 그름을 조사하고 그것에 그르게
의존한 일을 처벌할 강제관할권을 가지는 반면, 어떤 국제법원도 유사한
관할권을 가지지 않는다는 똑같이 명백한 사실에 비추어 평가되어야 한다.

이 의심스러운 유비들 중 일부는 국가들이 유엔 헌장 아래에서 부담한
의무들에 의해 훨씬 강화된 것으로 여겨질 수 있다. 그러나 여기서도, 그
강도에 대한 어떤 평가도, 종이 위에서는 훌륭한 헌장의 법집행
규정들(provisions)이 거부권 및 강대국들의 이념적 분열과 동맹들에 의해
마비되어 온 정도를 무시한다면 거의 가치가 없다. 때로 국내법의 법집행
규정들(provisions)도 총파업에 의해 마비될 \emphksans{수도} 있다고 답해지지만,
이것은 거의 설득력 있지 않다. 국내법과 국제법 사이의 비교에서 우리는
사실상 존재하는 것에 관심을 가지며, 여기서 사실들은 부인할 수 없이
다르기 때문이다.

그러나 국제법과 국내법 사이에 제안된 하나의 형식적 유비는 여기서 어느
정도 검토할 가치가 있다. 켈젠과 많은 현대 이론가들은 국내법처럼 국제법도
‘기초규범(basic norm)’, 또는 우리가 승인 규칙이라고 부른 것을 가지고
있고 기실 가져야만 한다고 고집한다. 체계의 다른 규칙들의 유효성은 그것을
참조하여 평가되며, 규칙들은 그것 덕분에 단일한 체계를 구성한다는 것이다.
반대 견해는 이러한 구조의 유비가 거짓이라는 것이다. 국제법은 단지 이러한
방식으로 통합되어 있지 않은, 분리된 의무의 일차 규칙들의
\emphksans{집합(set)}으로 이루어져 있다는 것이다. 국제법은 국제법학자들의
통상 용어로, 조약들에 구속력을 부여하는 규칙이 그중 하나인 관습 규칙들의
집합이다. 그 과제에 착수한 이들이 국제법의 ‘기초규범’을 정식화하는 데
매우 큰 어려움들을 발견했다는 것은 잘 알려져 있다. 이 지위의 후보들에는
\emph{pacta sunt servanda} 원칙이 포함된다. 그러나 이것은, 그 용어를
아무리 넓게 해석하더라도 국제법 아래의 모든 \textbf{(p.~234)} 의무들이
'\emph{pacta}'에서 생겨나는 것은 아니라는 사실과 양립하지 않는 것으로
보이기 때문에, 대부분의 이론가들에 의해 포기되었다. 그래서 그것은 덜
친숙한 어떤 것, 곧 이른바 ‘국가들은 자신들이 관습적으로 행위해 온 대로
행위해야 한다’는 규칙으로 대체되었다.

우리는 국제법의 기초규범에 관한 이러한 공식화들과 다른 경쟁적 공식화들의
이점들(merits)을 논의하지 않을 것이다. 그 대신 우리는 국제법이 그러한
요소를 반드시 포함해야 한다는 가정을 문제 삼을 것이다. 여기서 물어야 할
첫 질문이자 어쩌면 마지막 질문은 이것이다. 왜 우리는 이러한
\emphksans{선험적} 가정, 곧 그것은 바로 선험적 가정인데, 그것을 하고 그리하여
국제법 규칙들의 실제 성격을 선취 판단해야 하는가? 구성원들에게 의무들을
부과하는 규칙들을 ‘구속력 있는’ 것으로 삼아 살아가는 사회가, 그러한
규칙들을 어떤 더 기초적인 규칙에 의해 통합되거나 그 규칙에서 유효성을
도출하는 것이 아니라 단순히 분리된 규칙들의 집합으로 여길 수 있다는 것은
확실히 상상 가능하다. 어쩌면 그런 경우가 자주 있었을 수도 있다. 규칙들의
단순한 존재가 그러한 기초규칙의 존재를 수반하지 않는다는 것은 분명하다.
대부분의 현대 사회에는 예절 규칙들이 있고, 우리는 그것들이 의무들을
부과한다고 생각하지 않지만 그러한 규칙들이 존재한다고 말할 수 있다.
그러나 우리는 분리된 규칙들의 유효성이 도출될 수 있는 예절의 기초규칙을
찾지 않을 것이며, 찾을 수도 없을 것이다. 그러한 규칙들은 체계를 형성하지
않고 단순한 집합을 형성한다. 물론 예절보다 더 중요한 사안들이 걸려 있는
곳에서 이러한 사회통제 형식의 불편들은 상당하다. 그것들은 이미 제V장에서
서술되었다. 그러나 규칙들이 실제로 행위 기준들로 수용되고, 의무적
규칙들에 특징적인 적절한 형식의 사회적 압력에 의해 뒷받침된다면,
그것들이 구속력 있는 규칙들임을 보이기 위해 더 요구되는 것은 없다. 비록
이 단순한 사회구조 형식에서는 국내법에서 우리가 가지는 것, 곧 체계의
어떤 궁극적 규칙을 참조하여 개별 규칙들의 유효성을 논증하는 방식을
가지지 못하더라도 그러하다.

물론 체계가 아니라 단순한 집합을 구성하는 규칙들에 대해서도 우리가 물을
수 있는 여러 질문들이 있다. 예컨대 우리는 그것들의 역사적 기원에 관한
질문들이나 규칙들의 성장을 촉진한 인과적 영향들에 관한 질문들을 물을 수
있다. 또한 우리는 그것들에 따라 살아가는 이들에게 그 규칙들이 가지는
가치, 그리고 그들이 자신들을 그 규칙들에 복종하도록 도덕적으로 구속된
것으로 여기는지 \textbf{(p.~235)} 아니면 어떤 다른 동기에서 복종하는지를
물을 수 있다. 그러나 더 단순한 경우에는, 국내법이 그러하듯 기초규범 또는
이차적 승인 규칙에 의해 풍부해진 체계의 규칙들에 관하여 물을 수 있는 한
종류의 질문을 물을 수 없다. 더 단순한 경우 우리는 이렇게 물을 수 없다.
‘분리된 규칙들은 그 유효성 또는 “구속력(binding force)”을 체계의 어떤
궁극적 규정(provision)으로부터 도출하는가?’ 그러한 규정이 없고, 있을
필요도 없기 때문이다. 그러므로 기초규칙 또는 승인 규칙이 의무의 규칙들
또는 ‘구속력 있는’ 규칙들의 존재에 일반적으로 필요한 조건이라고 상정하는
것은 오류이다. 이것은 필수성이 아니라 사치이며, 그 구성원들이 분리된
규칙들을 하나씩 받아들이게 될 뿐 아니라 일반적 유효성 기준들에 의해
획정된 규칙들의 일반 부류들을 미리 수용하기로 결속되어 있는 발전된
사회체계들에서 발견된다. 더 단순한 사회 형식에서는 한 규칙이 규칙으로
수용되는지 아닌지를 기다려 보아야 한다. 기초적 승인 규칙을 가진
체계에서는 규칙이 실제로 만들어지기 전에, 그것이 승인 규칙의 요건들에
부합한다면 \emphksans{유효할 것}이라고 말할 수 있다.

같은 요점은 다른 형식으로 제시될 수 있다. 그러한 승인 규칙이 분리된
규칙들의 단순한 집합에 추가될 때, 그것은 체계의 이점들과 식별의 용이성을
함께 가져올 뿐 아니라, 처음으로 진술의 새로운 형식을 가능하게 한다.
이것들은 규칙들의 유효성에 관한 내부 진술들이다. 우리는 이제 새로운
의미에서 ‘체계의 어떤 규정(provision)이 이 규칙을 구속력 있게 만드는가?’
또는 켈젠의 언어로 ‘체계 안에서 그것의 유효성의 근거는 무엇인가?’라고
물을 수 있기 때문이다. 이 새로운 질문들에 대한 답들은 기초적 승인 규칙에
의해 제공된다. 그러나 더 단순한 구조에서 규칙들의 유효성이 어떤 더
기초적인 규칙을 참조하여 이렇게 논증될 수 없다고 하더라도, 이것은
규칙들이나 그 구속력 또는 유효성에 관하여 설명되지 않은 채 남겨지는 어떤
질문이 있다는 뜻이 아니다. 그러한 단순한 사회구조에서 규칙들이 왜 구속력
있는지에 관한 어떤 신비가 있고, 우리가 기초규칙을 찾을 수만 있다면
그것이 그 신비를 해소할 것이라는 말이 아니다. 단순한 구조의 규칙들은, 더
발전된 체계들의 기초규칙처럼, 수용되고 그러한 것으로 기능한다면 구속력
있다. 그러나 서로 다른 사회구조 형식들에 관한 이러한 단순한 진리들은,
이러한 바람직한 요소들이 사실상 발견되지 않는 곳에서 통일성과
\textbf{(p.~236)} 체계에 대한 완고한 탐색에 의해 쉽게 가려질 수 있다.

기실 기초규칙 없이 존재하는 가장 단순한 사회구조 형식들을 위해
기초규칙을 만들어 내려고 기울이는 노력에는 우스꽝스러운 점이 있다.
그것은 마치 벌거벗은 야만인이 \emphksans{정말로는} 어떤 보이지 않는 종류의
현대적 의복을 입고 있어야 \emphksans{한다고} 우리가 고집하는 것과 같다.
불행히도 여기에는 상존하는 혼동 가능성도 있다. 우리는 관련 사회,
개인들의 사회이든 국가들의 사회이든, 그것이 일정한 행위 기준들을 의무적
규칙들로 준수한다는 단순한 사실을 텅 빈 반복으로 말하는 어떤 것을
기초규칙으로 다루도록 설득될 수 있다. 국제법에 대해 제안되어 온 기묘한
기초규범, 곧 ‘국가들은 자신들이 관습적으로 행위해 온 대로 행위해야
한다’는 바로 이런 지위를 가지는 것이 분명하다. 그것은 일정한 규칙들을
수용하는 이들이 그 규칙들이 준수되어야 한다는 규칙도 준수해야 한다는 것
이상을 말하지 않기 때문이다. 이것은 규칙들의 집합이 국가들에 의해 구속력
있는 규칙들로 수용된다는 사실의 단순하고 무익한 중복일 뿐이다.

다시, 국제법은 기초규칙을 포함해야 \emphksans{한다}는 가정으로부터 우리가
해방되고 나면, 직면해야 할 질문은 사실의 질문이다. 국가들 사이의
관계들에서 기능하는 규칙들의 실제 성격은 무엇인가? 관찰될 현상들에 관한
서로 다른 해석들은 물론 가능하다. 그러나 여기서 제출되는 견해는,
국제법의 규칙들에 대한 일반적 유효성 기준들을 제공하는 기초규칙은 없고,
실제로 작동하는 규칙들은 체계가 아니라 규칙들의 집합을 구성하며, 그
가운데 조약들의 구속력을 마련하는 규칙들이 있다는 것이다. 많은 중요한
사안들에서 국가들 사이의 관계들이 다자조약들에 의해 규율되고, 이것들이
당사자가 아닌 국가들도 구속할 수 있다고 때때로 논변되는 것은 사실이다.
이것이 일반적으로 인정된다면, 그러한 조약들은 사실상 입법적 제정물들이
될 것이고 국제법은 그 규칙들에 대한 독자적인 유효성 기준들을 가지게 될
것이다. 그때에는 체계의 실제 특징을 나타내고 규칙들의 집합이 사실상
국가들에 의해 준수된다는 사실의 공허한 재진술 이상이 될 기초적 승인
규칙이 정식화될 수 있을 것이다. 아마 국제법은 현재 이것 및 구조상 그것을
국내체계에 더 가까워지게 할 다른 형식들의 수용으로 향하는 이행 단계에
있을지도 모른다. 만일, \textbf{(p.~237)} 그리고 이 이행이 완료될 때에는,
현재로서는 빈약하고 심지어 기만적으로 보이는 형식적 유비들이 실질을
획득할 것이며, 국제법의 법적 ‘성질(quality)’에 관한 회의론자의 마지막
의문들은 그때 잠재워질 수 있다. 그 단계에 도달할 때까지 유비들은 분명
형식의 유비들이 아니라 기능과 내용의 유비들이다. 기능의 유비들은 우리가
국제법이 도덕과 다른 방식들을 성찰할 때 가장 명료하게 나타나며, 그
일부를 우리는 앞 절에서 검토하였다. 내용의 유비들은 국내법과 국제법
모두에 공통되고 법률가들의 기법을 한쪽에서 다른 쪽으로 자유롭게 이전
가능하게 만드는 원칙들, 개념들, 방법들의 범위에 있다. ‘국제법’이라는
표현의 발명자인 벤담은 그것이 국내법과 ‘충분히 유비적’이라고 말하는
것만으로 그 표현을 방어하였다.\footnote{\emph{Principles of Morals and
  Legislation}, XVII. 25, n.~1.} 여기에 두 가지 언급을 덧붙일 가치가
있을 수 있다. 첫째, 그 유비는 형식의 유비가 아니라 내용의 유비라는 것,
둘째, 이 내용의 유비에서 국제법의 규칙들만큼 국내법에 가까운 다른 사회
규칙들은 없다는 것이다.

\newpage

\subsection{제10장 주석}\label{uxc81c10uxc7a5-uxc8fcuxc11d}


{\footnotesize
\noteentry 214쪽. \notetitle{‘국제법은 정말 법인가?’} --- 이것이 단지 사실문제로 오인된
언어적 질문에 불과하다는 견해는 Glanville Williams, 앞서 인용, \emph{22
BYBIL} (1945)을 참조하라.

\noteentry 215쪽. \notetitle{의문의 원천들(Sources of doubt).} --- 건설적인 일반적
개관으로는 A. H. Campbell, `International Law and the Student of
Jurisprudence' in \emph{35 Grotius Society Proceedings} (1950); Gihl,
`The Legal Character and Sources of International Law' in
\emph{Scandinavian Studies in Law} (1957)을 참조하라.

\noteentry 216쪽. \notetitle{‘국제법은 어떻게 구속력 있을 수 있는가?’} --- 이
질문(때로는 국제법의 구속력(binding force)의 문제로 지칭됨)은 Fischer
Williams, \emph{Chapters on Current International Law}, pp.~11--27;
Brierly, \emph{The Law of Nations}, 제5판 (1955), 제2장; \emph{The Basis
of Obligation in International Law} (1958), 제1장에서 제기된다. 또한
Fitzmaurice, `The Foundations of the Authority of International Law and
the Problem of Enforcement' in \emph{19 MLR} (1956)을 참조하라. 이
저자들은 규칙들의 체계가 구속력 있다(또는 구속력 있지 않다)는 단언의
의미를 명시적으로 논의하지 않는다.

\noteentry 217쪽. \notetitle{국제법에서의 제재(Sanctions).} --- 국제연맹 규약(League of
Nations Covenant) 제16조 하에서의 상황은 Fischer Williams, `Sanctions
under the Covenant' in \emph{17 BYBIL} (1936)을 참조하라. 유엔 헌장
제7장 하에서의 제재에 관해서는 Kelsen, `Sanctions in International Law
under the Charter of U.N.', \emph{31 Iowa LR} (1946), 그리고 Tucker,
`The Interpretation of War under present International Law', \emph{4 The
International Law Quarterly} (1951)을 참조하라. 한국전쟁에 대해서는
Stone, \emph{Legal Controls of International Conflict} (1954), 제9장,
담론 14(Discourse 14)를 참조하라. ‘평화를 위한 단결(유엔총회 결의)’이
유엔이 ‘마비되었다(paralyzed)’는 주장을 반박하는 것이라는 해석도
가능하다.

\noteentry 220쪽. \notetitle{의무적인 것으로 생각되고 말해지는 국제법.} --- Jessup,
\emph{A Modern Law of Nations}, 제1장, 그리고 `The Reality of
International Law', \emph{118 Foreign Affairs} (1940)을 참조하라.

\noteentry 220쪽. \notetitle{국가들의 주권(The Sovereignty of States).} --- ‘주권은 단지
국제 영역 중 법에 의해 국가들의 개별 행위에 남겨진 만큼에 부여된 이름일
뿐이다’라는 견해에 대한 명료한 설명은 Fischer Williams, 앞서 인용,
pp.~10--11, 285--99, \emph{Aspects of Modern International Law},
pp.~24--26, 그리고 Van Kleffens, `Sovereignty and International Law',
\emph{Recueil des Cours} (1953), I, pp.~82--83을 참조하라.

\noteentry 221쪽. \notetitle{국가(The State).} --- ‘국가(state)’라는 관념(notion)과
종속국가들(dependent states)의 유형에 대해서는 Brierly, \emph{The Law of
Nations}, 제4장을 참조하라.

\noteentry 224쪽. \notetitle{의사주의(voluntarist) 및 자기제한(auto-limitation) 이론들.} --- 주요 저자들로는 Jellinek, \emph{Die Rechtliche Natur der
Staatsverträge}; Triepel, `Les Rapports entre le droit interne et le
droit internationale', \emph{Recueil des Cours} (1923)이 있다. 극단적인
견해는 Zorn, \emph{Grundzüge des Völkerrechts}에 나타나며, 이에 대한
비판적 논의는 Gihl, 앞서 인용, \emph{Scandinavian Studies in Law}
(1957); Starke, \emph{An Introduction to International Law}, 제1장;
Fischer Williams, \emph{Chapters on Current International Law},
pp.~11--16을 참조하라.

\noteentry 224쪽. \notetitle{의무(obligation)와 동의(consent).} --- 어떤 국제법 규칙도
한 국가의 사전 동의, 명시적이든 묵시적이든 그러한 동의 없이는 그 국가에
구속력 있지 않다는 견해는 영국 법원들에서 표현되었으며(\emph{R.} v.
\emph{Keyn} 1876, 2 Ex. Div. 63, `The Franconia' 사건 참조),
상설국제사법재판소에서도 표현되었다. 이에 관해서는 \emph{The Lotus},
PCIJ Series A, No.~10을 참조하라.

\noteentry 226쪽. \notetitle{신생국 및 해양영역을 획득한 국가들.} --- 이에 대해서는
Kelsen, \emph{Principles of International Law}, pp.~312--13을 참조하라.

\noteentry 226쪽. \notetitle{비당사자들(non-parties)에 대한 일반 국제조약들의 효과.} ---
이에 대해서는 Kelsen, 앞서 인용, pp.~345 ff.; Starke, 앞서 인용, 제1장;
Brierly, 앞서 인용, 제7장, pp.~251--2를 참조하라.

\noteentry 227쪽. \notetitle{‘도덕(morality)’이라는 용어의 포괄적 사용.} --- 이에
대해서는 Austin, \emph{The Province} 제5강(Lecture V), pp.~125--9,
141--2에서 ‘실정 도덕(positive morality)’에 관한 논의를 참조하라.

\noteentry 230쪽. \notetitle{국제법에 복종할 도덕적 의무(moral obligation).} --- 이것이
국제법의 ‘토대(the foundation)’라고 보는 견해는 Lauterpacht, Brierly의
\emph{The Basis of Obligation in International Law} 서문(xviii), 그리고
Brierly, 같은 책 제1장을 참조하라.

\noteentry 232쪽. \notetitle{힘에 의해 부과된 조약을 입법으로 보는 것.} --- 이에
대해서는 Scott, `The Legal Nature of International Law', \emph{American
Journal of International Law} (1907), pp.~837, 862--4를 참조하라. 일반
조약들을 ‘국제 입법(international legislation)’이라고 부르는 통상적
서술에 대한 비판은 Jennings, `The Progressive Development of
International Law and its Codification', \emph{24 BYBIL} (1947),
p.~303에서 확인할 수 있다.

\noteentry 233쪽. \notetitle{탈중심화된 제재들(decentralized sanctions).} --- 이에
대해서는 Kelsen, 앞서 인용, p.~20; Tucker, 앞서 인용, \emph{4
International Law Quarterly} (1951)을 참조하라.

\noteentry 233쪽. \notetitle{국제법의 기초규범(basic norm).} --- 이를 \emph{pacta sunt
servanda}로 공식화한 논의는 Anzilotti, \emph{Corso di diritto
internazionale} (1923), p.~40을 참조하라. 이를 ‘국가들은 자신들이
관습적으로 행위해 온 대로 행위해야 한다’는 식으로 대체하려는 논의는
Kelsen, \emph{General Theory}, p.~369, 그리고 \emph{Principles of
International Law}, p.~418을 참조하라. 이에 대한 중요한 비판적 논의는
Gihl, \emph{International Legislation} (1937) 및 \emph{Scandinavian
Studies in Law} (1957), pp.~62 ff.에 있다. 국제법에 기초규범이 없다는
해석을 보다 충분히 전개한 논의로는 Ago, `Positive Law and International
Law', \emph{51 American Journal of International Law} (1957),
\emph{Scienza giuridica e diritto internazionale} (1958)을 참조하라.
Gihl은 국제사법재판소 규정 제38조에도 불구하고 국제법은 법의 형식적
원천들을 갖지 않는다고 결론짓는다. 국제법에 대해 하나의 ‘초기
가설(initial hypothesis)’을 정식화하려는 시도로서, 본문에서 제기된 것과
유사한 비판들에 열려 있는 것으로 보이는 견해는 Lauterpacht, \emph{The
Future of Law in the International Community}, pp.~420--3을 참조하라.

\noteentry 237쪽. \notetitle{국제법과 국내법 사이의 내용의 유비(analogy of content).} --- 이에 대해서는 Campbell, 앞서 인용, \emph{35 Grotius Society
Proceedings} (1950), p.~121 \emph{ad fin.}, 그리고 조약들, 영토취득,
시효, 임대, 위임통치, 지역권(servitudes) 등에 관한 규칙들을 논의한
Lauterpacht, \emph{Private Law Sources and Analogies of International
Law} (1927)을 참조하라.

\subsection{제10장 제3판
주석}\label{uxc81c10uxc7a5-uxc81c3uxd310-uxc8fcuxc11d}

\noteentry 220--226쪽. \notetitle{주권국가들과 국제법(Sovereign states and international law).} 주권(sovereignty)과 국제법의 제약들(constraints)의 양립가능성에
관해서는 Timothy Endicott, 「The Logic of Freedom and Power」, Samantha
Besson 및 John Tasioulas 편, \emph{The Philosophy of International Law}
(Oxford University Press, 2010)을 참조하라. 국가 주권(state
sovereignty)과 법의 지배(rule of law) 간의 관계에 대해서는 Jeremy
Waldron, 「The Rule of International Law」, (2006) \emph{Harvard Journal
of Law and Public Policy} 제30권, 15쪽을 참조하라. 초국가적
권위(transnational authority)의 다양한 등장 형태들에 비추어 주권을
재해석하려는 시도들에 대해서는 Neil MacCormick, \emph{Questioning
Sovereignty}를 참조하라. 현대 세계에서 ‘주권적’ 국가들의 사실상의
권력(de facto power)에 대한 의문들에 대해서는 Saskia Sassen,
\emph{Losing Control? Sovereignty in an Age of Globalization} (Columbia
University Press, 1996)을 참조하라.

\noteentry 227--232쪽. \notetitle{국제법과 도덕 (International law and morality).} 국제
영역에서 정의의 원칙들(principles of justice)의 적용에 대한 매우 영향력
있는 논의는 John Rawls, \emph{The Law of Peoples} (Harvard University
Press, 1999)을 참조하라. 또한 Allen Buchanan, \emph{Justice, Legitimacy,
and Self-Determination: Moral Foundations for International Law} (Oxford
University Press, 2003)을 참조하라. 국가들이 오직 자국의 국익을 증진하기
위해서만 국제법을 준수하고(comply with), 또 준수해야 한다는
‘현실주의(realist)’ 관점에 대해서는 Jack L. Goldsmith 및 Eric A. Posner,
\emph{The Limits of International Law} (Oxford University Press, 2005)를
참조하라.

\noteentry 232--237쪽. \notetitle{형식과 내용의 유비들 (Analogies of form and content).} Hart가 저술하던 시기에 비해 오늘날의 국제법은 더 체계화되었다는 시사에
대해서는 Samantha Besson 및 John Tasioulas가 편집한 \emph{The Philosophy
of International Law} 서문의 'Introduction'과 거기에서 인용된 원천들을
참조하라. 또한 Allen Buchanan 및 David Golove, 「Philosophy of
International Law」 in Jules L. Coleman 및 Scott Shapiro 편, \emph{The
Oxford Handbook of Jurisprudence and Philosophy of Law}를 참조하라.
}
